\chapter{Diskussion}
\label{chapter:kapitel5}
Einordnung und Interpretation der Evaluationsergebnisse im Gesamtkontext sowie kritische Reflexion der methodischen Grenzen, 
insbesondere der nachgelagerten Rollenübernahme als Entwickler nach Abschluss der Evaluation.

\section{Interpretation der Befunde im Gesamtkontext}
\label{gesamtkontext}

Von den sieben Interaktionsprinzipien der ISO~9241-110 sind vier vollständig, zwei
teilweise und ein Prinzip nicht erfüllt (vgl. Tabelle~\ref{tab:erfuellungsgrad}).
Über die prinzipienweise Betrachtung in Kapitel~\ref{section:ergebnisse} hinaus
lassen sich die elf Abweichungen zu drei Mustern verdichten.

Ein erstes Muster zeigt sich bei den Kriterien~2.1, 3.3, 6.3, 6.5, 6.6, 6.8 und~6.9.
Sie betreffen ausnahmslos Situationen, in denen die Anwendung auf Fehler, Ausnahmen
oder Rückmeldebedarf reagieren muss, etwa bei nicht erkennbaren aktiven Filtern,
fehlenden Warteinformationen, ungültigen Dateiformaten oder gescheiterten
Ladevorgängen. Die unmittelbar aufgabenbezogenen Prinzipien Aufgabenangemessenheit,
Erlernbarkeit und Steuerbarkeit sind dagegen vollständig erfüllt, was der in
Kapitel~\ref{section:problemstellung} beschriebenen Ausgangslage entspricht, wonach
die Anwendung bislang zwar funktional, jedoch nicht systematisch geprüft wurde.
Mit den Kriterien~6.5, 6.8 und~6.9 sind drei der vier Abweichungen
des Schweregrades~3 genau diesem Muster zuzuordnen.

Ein zweites Muster betrifft die sprachliche Verständlichkeit der Oberfläche. Die
Kriterien~2.4 und~3.4 (Meldungsart als SAP-Kürzel, Sprachmischung) sind mit
Schweregrad~3 beziehungsweise~2 eingestuft. Angesichts des in
Kapitel~\ref{section:untersuchungsgegenstand} beschriebenen Nutzungskontexts,
insbesondere der geringen SAP-Vertrautheit der rund 20 Nutzenden und der nur
zweimal wöchentlichen Nutzungsfrequenz, wiegen diese Befunde schwer. Das
erforderliche Vokabular prägt sich kaum ein, sodass die Mängel strukturell die
gesamte Zielgruppe betreffen.

Ein drittes, gegenläufiges Muster zeigt Kriterium~7.3 (fehlende Anpassung an
Handschuhnutzung). Trotz niedrigstem Schweregrad~1 weist es den unmittelbarsten
Bezug zu den physischen Rahmenbedingungen der Produktionshalle auf. Schweregrad
und Kontextrelevanz eines Befundes fallen demnach nicht notwendig zusammen.

Der aus Tabelle~A.1 der Norm entwickelte Kriterienkatalog
(Kapitel~\ref{section:untersuchungsdesign}) erweist sich damit, trotz der in
Kapitel~\ref{grenzen} diskutierten Einschränkungen, als hinreichend trennscharf.
Er macht sowohl die Erfüllung im Kernablauf als auch gezielte Abweichungen in
Randbereichen sichtbar. Für Woodward L'Orange knüpft dies unmittelbar an die in
Kapitel~\ref{section:problemstellung} beschriebene Ausgangslage an. An die Stelle
einer zuvor fehlenden objektivierbaren Grundlage tritt eine priorisierte Übersicht
identifizierter Schwachstellen (vgl. Tabelle~\ref{tab:optimierung}), wobei die
vier Abweichungen des Schweregrades~3 vorrangig zu adressieren
sind.\footcite[Vgl.][S.~47]{nielsen1994heuristic}

\section{Methodische Grenzen der Untersuchung}
\label{grenzen}

Eine methodische Besonderheit liegt in der Rollenkonstellation des Verfassers. An der
ursprünglichen Entwicklung der Anwendung war er nicht beteiligt und kannte diese zum
Zeitpunkt der Evaluation noch nicht (vgl. Kapitel~\ref{section:untersuchungsdesign}), sodass
das für Inspektionsverfahren relevante Risiko einer Betriebsblindheit gegenüber selbst
getroffenen Gestaltungsentscheidungen zum Evaluationszeitpunkt nicht bestand. Eine
Einschränkung ergibt sich jedoch aus der nachgelagerten Rollenübernahme. Da der Verfasser die
identifizierten Abweichungen im Anschluss selbst als Entwickler beheben wird, lässt sich nicht
ausschließen, dass diese bevorstehende Zuständigkeit die Schweregradbeurteilung unbewusst
beeinflusst hat, etwa indem schwerer zu behebende Probleme niedriger eingestuft werden, um den
eigenen späteren Aufwand zu begrenzen. Eine Kalibrierung durch einen zweiten, an der Behebung
unbeteiligten Gutachter fand nicht statt.

Diese Einschränkung verschärft sich dadurch, dass nur ein Gutachter tätig war, während für die
heuristische Evaluation üblicherweise drei bis fünf Gutachter empfohlen
werden.\footcite[Vgl.][S.~254 f.]{nielsen1990heuristic} Nielsen und Landauer zeigen anhand von elf
Fallstudien, dass sich das Auffinden von Usability-Problemen mit zunehmender Gutachterzahl als
Poisson-Prozess beschreiben lässt und ein einzelner Gutachter im Mittel nur etwa ein Drittel
der insgesamt vorhandenen Probleme
findet.\footcite[Vgl.][S.~206--207]{nielsen1993mathematical} Entscheidender ist eine
strukturelle Eigenschaft ihres Modells. Da sich die Anzahl gefundener Probleme aus zwei
zunächst unbekannten Größen zusammensetzt, lassen sich Gesamtzahl und Fundwahrscheinlichkeit
erst ab mindestens zwei unabhängigen Evaluationsdurchläufen
schätzen.\footcite[Vgl.][S.~208 f.]{nielsen1993mathematical} Bei nur einem Gutachter lässt sich
folglich nicht einmal die Größenordnung der unentdeckten Probleme schätzen. Die elf
dokumentierten Abweichungen sind daher als Teilmenge, nicht als vollständige Erhebung zu
verstehen.

Eine weitere Grenze liegt im Verfahren selbst. Die heuristische Evaluation prognostiziert
Usability-Probleme, verifiziert deren tatsächliches Auftreten im realen Arbeitsablauf jedoch
nicht.\footcite[Vgl.][S.~254 f.]{nielsen1990heuristic} Ob die Befunde tatsächlich zu den in
Kapitel~\ref{section:problemstellung} beschriebenen Bedienfehlern und Zeitverlusten führen,
bleibt daher offen. Hinzu kommt, dass der in Kapitel~\ref{section:untersuchungsgegenstand}
beschriebene Nutzungskontext wesentlich auf einem Gespräch mit dem Meister der Instandhaltung
beruht und nicht auf eigener Beobachtung unter Realbedingungen.

Schließlich schränken zwei weitere Faktoren die Aussagekraft ein. Der Ausschluss nicht
anwendbarer Norm-Empfehlungen (Kapitel~\ref{section:untersuchungsdesign}) beruht auf einer
Einschätzung des Evaluators, und die Untersuchung bildet eine Momentaufnahme eines bestimmten
Versionsstands, dessen Altsystemkontext aus SAP~R/3 mit SAP Gateway den Umsetzungsspielraum der
Optimierungsvorschläge zusätzlich begrenzt (vgl. Kapitel~\ref{section:optimierungspotenziale}).